\section{Sum of Subarray Ranges}
\label{sec:sq-sum-subarray-ranges}

\problemstatement{Given an array $\textit{arr}$, the \textbf{range} of a
subarray is defined as its maximum element minus its minimum element.
Compute the sum of ranges over all contiguous subarrays.}

\begin{patternbox}
A subarray's range is, by definition, $\max - \min$. Summing ranges over
all subarrays is therefore exactly
\[
  \left(\sum_{\text{subarrays}} \max\right) -
  \left(\sum_{\text{subarrays}} \min\right),
\]
and the second term is precisely Section~\ref{sec:sq-sum-subarray-minimums}'s
problem. The first term is its exact mirror image -- ``sum of subarray
maximums'' -- solved by the same contribution-counting technique with
every comparison flipped. No new derivation is needed; this problem is a
direct combination of two near-identical applications of the same idea.
\end{patternbox}

\subsection*{Understanding the problem carefully}

Linearity of summation lets us split the sum of $(\max - \min)$ over all
subarrays into a sum of $\max$ over all subarrays, minus a sum of
$\min$ over all subarrays, computed completely independently. Each piece
is solved by the contribution-counting technique from
Section~\ref{sec:sq-sum-subarray-minimums}: for the minimum part, exactly
as before; for the maximum part, the same left/right boundary counting,
but searching for the nearest strictly-\emph{greater} element (instead of
strictly-less) on one side and greater-or-equal on the other.

\begin{ideabox}[Reuse Twice, with Comparisons Flipped for the Maximum Half]
Compute $\textit{sumOfMins}$ exactly as in
Algorithm~\ref{sec:sq-sum-subarray-minimums}. Compute
$\textit{sumOfMaxs}$ with the same algorithm, but with stack comparisons
flipped: the left pass pops while the stack's top value $\le
\textit{arr}[i]$ (finding the nearest strictly \emph{greater} element to
the left), and the right pass pops while the stack's top value $<
\textit{arr}[i]$ (finding the nearest greater-or-equal element to the
right). Return $\textit{sumOfMaxs} - \textit{sumOfMins}$.
\end{ideabox}

\subsection*{Why flipping every comparison correctly computes the maximum-side sum}

This is a direct restatement of the next-greater/next-smaller duality
established in Section~\ref{sec:sq-next-smaller}: ``how many subarrays is
$\textit{arr}[i]$ the maximum of'' is structurally identical to ``how many
subarrays is $\textit{arr}[i]$ the minimum of,'' with every $<$ replaced
by $>$ and vice versa. The strict/non-strict asymmetry between the two
boundary searches is still required, for the same reason as
Section~\ref{sec:sq-sum-subarray-minimums}: to avoid double-counting
subarrays containing duplicate maximum values.

\subsection*{Algorithm}

\begin{enumerate}
  \item Compute $\textit{sumOfMins}$ using
        Algorithm~\ref{sec:sq-sum-subarray-minimums} unchanged.
  \item Compute $\textit{sumOfMaxs}$ using the same algorithm structure,
        with left-pass pops on $\le$ and right-pass pops on $<$ (instead
        of $\ge$ and $>$ respectively).
  \item Return $\textit{sumOfMaxs} - \textit{sumOfMins}$.
\end{enumerate}

\subsection*{Dry run}

\dryrun{Take $\textit{arr} = [1,2,3]$.}

$\textit{sumOfMins}$: subarrays and their minimums: $[1]{=}1$, $[2]{=}2$,
$[3]{=}3$, $[1,2]{=}1$, $[2,3]{=}2$, $[1,2,3]{=}1$. Sum $=1+2+3+1+2+1=10$.

$\textit{sumOfMaxs}$: maximums: $1,2,3,2,3,3$. Sum
$=1+2+3+2+3+3=14$.

Sum of ranges $=14-10=4$. Checking directly: ranges are $0,0,0,1,1,2$,
summing to $4$. Matches.

\subsection*{C++ implementation}

\begin{lstlisting}
#include <bits/stdc++.h>
using namespace std;

long long sumOfMins(vector<int>& arr) {
    int n = arr.size();
    vector<int> left(n), right(n);
    stack<int> st;

    for (int i = 0; i < n; i++) {
        while (!st.empty() && arr[st.top()] >= arr[i]) st.pop();
        left[i] = i - (st.empty() ? -1 : st.top());
        st.push(i);
    }
    while (!st.empty()) st.pop();
    for (int i = n - 1; i >= 0; i--) {
        while (!st.empty() && arr[st.top()] > arr[i]) st.pop();
        right[i] = (st.empty() ? n : st.top()) - i;
        st.push(i);
    }

    long long total = 0;
    for (int i = 0; i < n; i++) total += (long long)arr[i] * left[i] * right[i];
    return total;
}

long long sumOfMaxs(vector<int>& arr) {
    int n = arr.size();
    vector<int> left(n), right(n);
    stack<int> st;

    for (int i = 0; i < n; i++) {
        while (!st.empty() && arr[st.top()] <= arr[i]) st.pop();
        left[i] = i - (st.empty() ? -1 : st.top());
        st.push(i);
    }
    while (!st.empty()) st.pop();
    for (int i = n - 1; i >= 0; i--) {
        while (!st.empty() && arr[st.top()] < arr[i]) st.pop();
        right[i] = (st.empty() ? n : st.top()) - i;
        st.push(i);
    }

    long long total = 0;
    for (int i = 0; i < n; i++) total += (long long)arr[i] * left[i] * right[i];
    return total;
}

long long sumSubarrayRanges(vector<int>& arr) {
    return sumOfMaxs(arr) - sumOfMins(arr);
}

int main() {
    vector<int> arr = {1, 2, 3};
    cout << "Sum of subarray ranges: " << sumSubarrayRanges(arr) << endl;
    return 0;
}
\end{lstlisting}

\begin{complexitybox}
\textbf{Time Complexity.} Four monotonic-stack passes total (two for
minimums, two for maximums), each $O(n)$, giving
\[
  O(n).
\]
\textbf{Space Complexity.} $O(n)$.
\end{complexitybox}

\begin{takeawaybox}
When a quantity decomposes linearly (here, range $= \max - \min$),
compute each piece with whatever technique already solves it, and combine
at the end -- there is no need to find a single unified algorithm that
handles $\max$ and $\min$ simultaneously when summing them separately and
subtracting is just as correct and far simpler to reason about.
\end{takeawaybox}
